%
% komplexitaet.tex -- Beispiel-File für Teil 3
%
% (c) 2026 Nathanael Fässler, OST Ostschweizer Fachhochschule
%
% !TEX root = ../../buch.tex
% !TEX encoding = UTF-8
%

\section{Komplexität}
Dieser Abschnitt setzt theoretische Grundlagen zur \emph{Komplexitätstheorie} voraus, die sehr gut von Müller in seinem Buch~\cite{buchberger:automaten-und-sprachen} erklärt werden.
Bekanntlich ist das Lösen eines Sudokus ein NP-vollständiges Problem.
Die Klasse \emph{NP} beinhaltet alle entscheidbaren Sprachen, die sich von einer nichtdeterministischen Maschine\footnote{Leider existieren nichtdeterministische Maschinen nur in der Theorie, 
und daher lassen sich diese Probleme in der Praxis nur in exponentieller Zeit lösen.} 
in polynomieller Zeit lösen lassen.
Als \emph{NP-vollständig} bezeichnet man die Klasse der schwierigsten Probleme innerhalb von NP.
Es stellt sich die Frage, ob das Lösen von Polynomgleichungssystemen auch NP-vollständig ist.

\subsection{Entscheidungsprobleme}
Formell lässt sich Sudoku als das Entscheidungsproblem
\begin{equation}
  \begin{split}
    \text{SUDOKU} =
    \left\{
      \langle\mathcal{S}\rangle
      \;\left|\;
        \begin{minipage}{7cm}\raggedright
          $\mathcal{S}$ ist ein Sudoku mit vorgegebenen Zahlen, welches eine Lösung besitzt, welche die bekannten Sudoku-Regeln erfüllt.
        \end{minipage}
      \right.
    \right\}
  \end{split}
  \label{buchberger:sudoku-problem}
\end{equation}
definieren.
Das Entscheidungsproblem beantwortet nur die Frage, ob es eine Lösung für das Sudoku gibt oder nicht. Es findet nicht die Lösung selbst.
Es wurde jedoch bewiesen, dass das Entscheiden der Existenz einer Lösung in der gleichen Komplexitätsklasse liegt wie das Finden der Lösung.

Für das Problem der Polynomgleichungssysteme wird das Entscheidungsproblem
\begin{equation}
  \begin{split}
    \text{POLYNOM} =
    \left\{
      \langle\mathcal{P}\rangle
      \;\left|\;
        \begin{minipage}{7cm}\raggedright
          $\mathcal{P}$ ist ein System von Polynomgleichungen über dem endlichen Körper $\mathbb{F}_p$, welches mindestens eine Lösung besitzt.
        \end{minipage}
      \right.
    \right\}
  \end{split}
  \label{buchberger:polynom-problem}
\end{equation}
definiert.
Die Einschränkung auf Polynomgleichungen über dem endlichen Körper $\mathbb{F}_p$ ist hier sinnvoll, da der Körper ausreicht um die Zahlen eines Sudokus zu codieren.
Für unendliche Zahlenbereiche ist es komplizierter, wie das zehnte hilbertsche Problem zeigt.

\subsection{Ist $POLYNOM$ NP-vollständig?}
\begin{theorem}
  $POLYNOM$ ist NP-vollständig.
\end{theorem}

\subsubsection{Reduktion}
Wie gezeigt, lässt sich $SUDOKU$ durch $POLYNOM$ lösen.
Somit ist $POLYNOM$ mindestens so schwierig wie $SUDOKU$.
Dieser Sachverhalt lässt sich durch die Reduktion
$SUDOKU \leq_{\mathrm{p}} POLYNOM$
ausdrücken.

Dies genügt nicht, um zu beweisen, dass $POLYNOM$ NP-vollständig ist.
Es könnte ja sein, dass $POLYNOM$ noch schwieriger ist als NP.

\subsubsection{Entscheidbarkeit}
Ein Entscheidungsproblem gilt als entscheidbar, wenn es möglich ist, das Problem in endlicher Zeit zu entscheiden.
Da $POLYNOM$ auf den endlichen Körper $\mathbb{F}_p$ eingeschränkt ist, gibt es nur endlich viele Kombinationen aller Variablen und Zahlen, welche einfach durchprobiert werden können.
Somit ist $POLYNOM$ entscheidbar.

\subsubsection{NP-Zugehörigkeit}

Das Problem $POLYNOM$ lässt sich in polynomieller Zeit verifizieren.
Das heisst es gibt einen Algorithmus, mit dem man in polynomieller Zeit prüfen kann ob eine vorgeschlagene Lösung tatsächlich eine Lösung des Gleichungssystems ist.

Folgender Algorithmus ist ein solcher Verifizierer:
\begin{enumerate}
  \item
  Der Algorithmus erhält ein Gleichungssystem und eine Lösung in der Form $x_1 = 5, \dots x_n = 12$.
  \item
  Setze in jeder Gleichung für jede Variable den in der Lösung angegebenen Wert ein.
  \item
  Rechne in jeder Gleichung die Zahlen aus.
  \item
  Überprüfe für jede Gleichung ob der Wert auf der linken und rechten Seite identisch ist.
\end{enumerate}
Die Laufzeit hängt von den Anzahl Gleichungen $n$, den Anzahl Variablen $v$ und der höchsten Potenz $p$ ab.
Schritt 2 besitzt die Laufzeit $O(n \cdot v)$.
Das Ausrechnen der Zahlen in Schritt 3 hat Laufzeit $O(n \cdot v \cdot p)$. Das Ausrechnen einer Potenz $x^p$ benötigt $p$ Multiplikationen, welche polynomielle Laufzeit haben.
Der letzte Schritt besitzt die Laufzeit $O(n)$.

Somit ergibt sich die Gesamtlaufzeit
\begin{equation}
O(n \cdot v) + O(n \cdot v \cdot p) + O(n) = O(n \cdot v \cdot p)
,
\end{equation}
welche polynomiell ist.

Zusammen mit der zuvor gezeigten Entscheidbarkeit zeigt dieser Algorithmus, dass $POLYNOM$ in NP liegt und nicht ausserhalb.

\begin{proof}
Die NP-vollständigen Probleme sind die Klasse der schwierigsten Probleme innerhalb von NP und sind alle gleich schwierig.
Durch die beiden Resultate
\begin{enumerate}
  \item
  $POLYNOM$ ist mindestens so schwierig wie $SUDOKU$.
  \item
  $POLYNOM$ liegt in NP
\end{enumerate}
lässt sich schlussfolgern, dass $POLYNOM$ gleich schwierig ist wie $SUDOKU$.
Dies lässt sich durch
$SUDOKU \equiv_{\mathrm{P}} POLYNOM$
ausdrücken.
Somit ist bewiesen, dass $POLYNOM$ NP-vollständig ist.
\end{proof}

\subsubsection{Polynomgleichungssysteme über den ganzen Zahlen}
Die Frage zur NP-vollständigkeit grösserer Zahlenbereiche als $\mathbb{F}_p$ ist kompliziert.
Zum Beispiel bei einem Polynomgleichungssystem über $\mathbb{Z}$ ist die Frage, ob eine Lösung existiert, nicht entscheidbar.
Dies war der Inhalt des 10. Hilbertschen Problems, welches fragte, ob es möglich ist für jede Polynomgleichung zu entscheiden, ob es eine Lösung gibt.
1970 wurde das Problem gelöst, mit der Antwort, dass es keinen solchen Entscheider geben kann.
Somit ist das Entscheidungsproblem der Polynomgleichungssysteme über $\mathbb{Z}$ ausserhalb von NP, also schwieriger, und somit nicht NP-vollständig.

\subsubsection{Polynomgleichungssysteme über den reellen Zahlen}
Bei Polynomgleichungssystemen über $\mathbb{R}$ ist die Frage, ob es eine Lösung gibt, erstaunlicherweise entscheidbar, wie Tarski bewies.
Allerdings ist das Problem nicht verifizierbar, aufgrund der unendlich langen Dezimalbrüche.
Ein Verifizierer bräuchte unendlich lange Rechenzeit schon nur um eine irrationale Zahl einzulesen.


\subsection{P vs. NP}
Die Millennium-Probleme umfassen eine Liste von ungelösten mathematischen Problemen.
Wer eines davon löst wird mit einer Million US-Dollar belohnt.
Eines der Millennium-Probleme ist die Frage $\text{P} \stackrel{?}{=} \text{NP}$, also ob die Komplexitätsklassen P und NP verschieden sind, was von den meisten Mathematikern vermutet wird.
Falls sich herausstellen würde, dass $\text{P} = \text{NP}$ gilt, würde es bedeuten, dass es für alle NP-vollständigen Probleme einen Algorithmus gibt, der das Problem in polynomieller Zeit löst.
Wenn also jemand es schafft, einen allgemeinen Algorithmus zu finden der Sudokus oder Polynomgleichungssysteme, welche wie besprochen NP-vollständig sind, löst, hätte $\text{P} = \text{NP}$ bewiesen.
Dies hätte schwerwiegende Auswirkungen zum Beispiel auf die Kryptographie:
Die meisten kryptographischen Systeme basieren auf einer langen Rechenzeit um einen unbekannten \emph{Schlüssel} zu berechnen.
Der beste Weg so einen Key zu berechnen ist meistens durch Bruteforcing\footnote{Alle möglichen Varianten ausprobieren.}.
Wenn es auf einmal einen Weg gibt Keys in polynomieller Zeit zu berechnen, wären viele kryptographische Systeme nicht mehr sicher.
